% Separate post-CP28 candidate; no claim about larger or reducible wild images.
\subsection{Excluding the actual fourth-order Cartan comparison}

Throughout this subsection, \(K\) is complete discretely valued of
mixed characteristic \((0,p)\), its residue field \(k\) is
algebraically closed, \(p\ge5\), \(e_K=p(p-1)\), and
\(\mu_p\subset K\). All comparisons concern actual generic
\(p^3\)-torsion of potentially supersingular height-two groups.
Put \(q=p^2\), \(N=q-1\), and use the covariant special-character
convention of B53.

\begin{lemma}[The two actual Cartan displacement levels]
\label{lem:B-RD81-levels}
\status{Candidate; full proof supplied, independent audits pending.}
Suppose that two such groups have an actual \(p^3\)-torsion
isomorphism, irreducible residual inertia, and distinct finite
special characters. If their common wild third-level image has
order \(p^4\), use the common fields
\[
 L/K\text{ of tame degree }N,\qquad
 F=L(M_1[p^3])=L(M_2[p^3]).
\]
Then \(W=\operatorname{Gal}(F/L)\simeq(\mathbf Z/p^2)^2\),
\(U=pW\) is its terminal kernel of order \(q\), and both actual
matrix layers are the original nonsplit Cartan. For
\[
 A=p(p-1),\qquad D_1=qA,\qquad D_2=q^2A,
 \qquad v_F(p)=ND_2,                                  \tag{B-RD81.1}
\]
each primitive coordinate has valuation \(A\); its displacement
by \(g\in W\setminus U\) has valuation \(D_1\), and by
\(h\in U\setminus\{1\}\) has valuation \(D_2\).

Let \(a\) be the minimum lower break of nonidentity elements
of \(W\), with the convention
\(b(g)=v_F(g\varpi-\varpi)-1\). Every nonidentity element
of \(U\) has break strictly greater than \(a\), so a
minimum-break element lies outside \(U\).
\end{lemma}
\begin{proof}
The actual Cartan conclusion, including \(W=(\mathbf Z/p^2)^2\)
and \(pW=U\), is Theorem~\ref{thm:B-RD80}. The valuation
statements are (B-RD80.1). A nonzero Cartan coefficient is an
invertible multiplication operator on the residual two-dimensional
space. Hence a nonterminal element moves every primitive point
by a point of exact order \(p^2\), and a nonidentity terminal
element moves it by a nonzero point of order \(p\).

For completeness the strict-break assertion uses the actual
group, not its associated graded dimensions. If \(g\) has
break \(b\), the coefficient
\[
 d_g=\overline{(g\varpi-\varpi)/\varpi^{b+1}}
\]
adds under composition of elements of break at least \(b\).
Every such element fixes the residue field, so its \(p\)th
power has zero coefficient in that degree. Thus
\(b(g^p)>b(g)\) whenever \(g^p\ne1\). Every nonidentity
\(h\in U=pW\) is the \(p\)th power of an order-\(p^2\)
element of \(W\), which proves \(b(h)>a\).
\end{proof}

\begin{proposition}[The full finite Cartan does not conceal a different character]
\label{prop:B-RD81-Cartan}
\status{Candidate; full proof supplied, independent audits pending.}
Under the field hypotheses above, an actual \(p^3\)-congruent
pair with irreducible residual inertia and wild image of order
\(p^4\) has equal finite special characters.
\end{proposition}
\begin{proof}
Suppose the characters differ and apply the preceding lemma.
Fix one uniformizer \(\varpi\) of the actual common field \(F\).
For each model fix Honda coordinates and, once and for all,
primitive point lifts indexed by \(z\in\mathbf F_q^\times\).
Their leading coordinates are \(d_i z\), where \(d_i\in k^\times\).
The coordinate expansions are taken using residue representatives
in a coefficient ring of \(K\), and are written
\[
 x_i(z)=\sum_{r\ge A}[f_{i,r}(z)]\varpi^r,
 \qquad f_{i,A}(z)=d_i z.                             \tag{B-RD81.2}
\]
The representatives are fixed by \(W\). Since \(v_F(p)>D_2\),
all terms relevant up to displacement valuation \(D_2\) are
calculated in characteristic \(p\), exactly as in
Lemma~\ref{lem:B-RD80-first}. These are fixed point lifts
throughout the argument; no independence of higher coefficients
from the choice of lift is asserted.

Write \(\lambda_i(g)\in\mathbf F_q^\times\) for the first
matrix coefficient of \(g\in W\setminus U\), and
\(\lambda_i(h)\in\mathbf F_q^\times\) for the terminal
coefficient of \(h\in U\setminus\{1\}\). The corresponding
coefficient maps on \(W/U\) and \(U\) are additive isomorphisms
onto \(\mathbf F_q\). The actual leading displacement
coefficients in the two levels are
\[
 d_i^q\lambda_i(g)z,
 \qquad d_i^{q^2}\lambda_i(h)z.                       \tag{B-RD81.3}
\]
Indeed the leading \(p\)-series at a primitive point is
\(T^q\), uniquely, and multiplication by \(p^2\) raises
leading coordinates to the \(q^2\)th power. Normalized field
coordinates satisfy \(z^q=z\). Since the two special
characters differ, the actual intertwiner has Frobenius
orientation on these leading spaces, as proved in B79.
Consequently, on both matrix layers,
\[
                   \lambda_2=\lambda_1^p.             \tag{B-RD81.4}
\]
The possible scalar in the intertwiner cancels in conjugating
Cartan multiplication operators.

Set
\[
 a_0=\frac{D_1-A}{p}=(q-1)(p-1),\qquad
 j_0=A+(p-1)a_0=D_1-a_0.                             \tag{B-RD81.5}
\]
Here \(a_0\equiv1\pmod p\), and hence
\(j_0\equiv-1\pmod p\). In particular the boundary degree
\(j_0\) is never divisible by \(p\).
Choose a minimum-break element \(g\in W\setminus U\), of
break \(a\). The degree-\(A\) term of (B-RD81.2) first
changes at degree \(A+pa\), with coefficient
\((A/p)d_i z d_g^p\). Any other degree divisible by \(p\)
changes no earlier than \(A+pa\). The first nonzero degree
\(j\) prime to \(p\) changes at \(j+a\), with coefficient
\(j f_{i,j}(z)d_g\); it is the unique earliest prime-degree
change.

First suppose \(a<a_0\). Cancellation below \(D_1\)
forces the same first prime degree in both models,
\[
                  j=A+(p-1)a,
          \qquad f_{i,j}(z)=c_i z,
          \qquad c_i\ne0.                            \tag{B-RD81.6}
\]
This is the low-break conclusion of
Lemma~\ref{lem:B-RD80-first}; it also forces \(p\nmid a\).
If instead \(a>a_0\), the first change at \(D_1\) has to
come from a prime degree, and (B-RD81.3) gives
\[
                 j=D_1-a,
          \qquad f_{i,j}(z)=c_i z,
          \qquad c_i\ne0.                            \tag{B-RD81.7}
\]
If that degree is not a positive permitted prime degree, the
case is already impossible.

In either case every nonidentity \(h\in U\) has break
\(b>a\), and this fixed first prime term is uniquely the
first term changed by \(h\). To check the strict comparison,
the difference \(A+pb-(j+b)\) equals respectively
\[
 (p-1)(b-a),\qquad
 (p-1)(b-a)+p(a-a_0),                                \tag{B-RD81.8}
\]
both positive. All later prime terms change later as well.
Thus \(j+b=D_2\) for every nonidentity terminal element.
Its leading coefficient yields
\[
      d_i^{q^2}\lambda_i(h)z=j c_i z d_h,
      \qquad \lambda_i(h)=\kappa_i d_h,
      \qquad \kappa_i\ne0.                           \tag{B-RD81.9}
\]
The common lower break makes \(d_h\) the same actual
uniformizer coefficient for the two models. Since
\(\lambda_1(U)=\mathbf F_q\), equations (B-RD81.4) and
(B-RD81.9) would make Frobenius on \(\mathbf F_q\) a fixed
scalar map. Evaluating at \(1\) and at an element outside
\(\mathbf F_p\) is a contradiction. Therefore \(a=a_0\).

We now treat the boundary, including a vanishing first prime
coefficient. No prime degree below \(j_0\) can occur in
either expansion: its change would precede \(D_1\), whereas
all divisible degrees first change at least at \(D_1\).
At degree \(D_1\) the only possible contributions are
the initial term and the degree-\(j_0\) term. Thus, for
the fixed minimum-break element,
\[
 d_i^q\lambda_i(g)z
   =(A/p)d_i z d_g^p+j_0 f_{i,j_0}(z)d_g.             \tag{B-RD81.10}
\]
Since \(j_0d_g\ne0\), this proves
\(f_{i,j_0}(z)=c_i z\), with \(c_i\) possibly zero.

If \(c_i\ne0\) for at least one model, then for every
terminal element its break \(b>a_0\) makes this term
uniquely first, because
\(A+pb-(j_0+b)=(p-1)(b-a_0)>0\). Hence every such
break is \(b_*=D_2-j_0\). In the other model the same
actual terminal element must also have a nonzero first
prime coefficient at degree \(j_0\): its divisible terms
all change after \(D_2\), and an earlier or later first
prime term would give a displacement valuation different
from \(D_2\). Equivalently,
\[
 b_*-\frac{D_2-A}{p}
       =\frac{p-1}{p}(D_2-D_1)>0,                    \tag{B-RD81.11}
\]
so the high-break clause of Lemma~\ref{lem:B-RD80-first}
forces that exact nonzero degree. We obtain (B-RD81.9)
for both models and the same contradiction. It follows
that \(c_1=c_2=0\).

There is consequently no prime degree at or below \(j_0\)
in either expansion. No nonterminal element can now have
break \(b>a_0\): all divisible terms would change after
\(A+pa_0=D_1\), and all prime terms after
\(j_0+a_0=D_1\). This contradicts its actual displacement
valuation \(D_1\). Every element of \(W\setminus U\)
therefore has break \(a_0\), and every nonidentity element
of \(U\) has greater break. The leading ramification map
is an injective additive map \(W/U\to k\) with image of
dimension two over \(\mathbf F_p\).

For all nonterminal elements equation (B-RD81.10) now has
only its first term. Consequently
\[
       \lambda_i(g)=(A/p)d_i^{1-q}d_g^p
                   \quad(g\in W\setminus U).          \tag{B-RD81.12}
\]
The two coefficient maps are again fixed scalar multiples
of one another, while (B-RD81.4) makes one the Frobenius
of the other. Their values fill \(\mathbf F_q\), so this
is impossible. The contradiction proves the proposition.
\end{proof}

\begin{corollary}[The first two possible wild orders are excluded]
\label{cor:B-RD81}
\status{Candidate; full proof supplied, independent audits pending.}
At \(e=p(p-1)\), an actual \(p^3\)-congruent pair with
irreducible residual inertia and different finite special
characters must have wild third-level image of order \(p^r\)
with \(5\le r\le8\). In particular an irreducible-residual
counterexample at \(p=5,e=20\) would have actual
\(125\)-torsion wild image of order at least \(3125\).
\end{corollary}
\begin{proof}
B79 excludes order \(p^3\), and B80 together with the
preceding proposition excludes order \(p^4\). The upper
bound \(p^8\) is the actual matrix-kernel bound of B79.
\end{proof}

\begin{remark}[Exact scope of the Cartan exclusion]
\label{rem:B-RD81-scope}
\status{Candidate scope; the larger and reducible images remain open.}
The argument excludes an actual finite-level Cartan comparison
using its point coordinates and its exact \(p\)th-power map.
It does not promote a finite-level quadratic action to a
finite-extension endomorphism of the whole group and does
not invoke the potentially CM theorem for a non-CM model.
Wild image orders \(p^5,\ldots,p^8\) and reducible residual
comparisons remain open. Thus the full first-wild least
index still lies in \(\{2,3\}\), with no counterexample or
unrestricted \(125\)-recovery assertion.
\end{remark}
